% !PRINTED 598
\weberpage{598}
Die Zahl
\weq{4}{h = (\overline{O},\, \overline{E}A)}
heißt die \gesperrt{Klassenzahl des Körpers $\Omega$ nach $A$}, und das System
$A\mathfrak{a}_1, A\mathfrak{a}_2, \ldots, A\mathfrak{a}_h$ die \gesperrt{Idealklassen nach $A$}. Den Satz 2. können
wir nach den allgemeinen Gruppensätzen im \S\,100, 12. beweisen.
Danach ist:
\[
\begin{aligned}
(\overline{O},\, \overline{E}A) &= (\overline{O},\, \overline{E}O)\,(\overline{E}O,\, \overline{E}A), \\
(\overline{E}O,\, \overline{E}A) &= (\overline{E}O,\, \overline{E}EA) = (O,\, EA), \\
(O,\, EA)\,(EA,\, A) &= (O,\, A), \\
(EA,\, A) &= (E,\, E'),
\end{aligned}
\]
wenn $E'$ die Gruppe der in $A$ enthaltenen numerischen Einheiten
bedeutet. Folglich ist
\weq{5}{(\overline{O},\, \overline{E}A) = (\overline{O},\, \overline{E}O)\,\dfrac{(O,\, A)}{(E,\, E')},}
und $(\overline{O},\, \overline{E}O)$ ist die Zahl der absoluten Idealklassen in $\Omega$, also
eine endliche Zahl. Die Zahl $(E,\, E')$ ist im quadratischen Körper
mit negativer Diskriminante gleich 1, wenn $-1$ in $A$ vorkommt,
gleich 2, wenn $-1$ nicht in $A$ vorkommt, und kann nur in den
beiden Ausnahmefällen $\Delta = -4$, $\Delta = -3$ bis zu 4 oder bis
zu 6 ansteigen.

Gehören $\mathfrak{a}_1, \mathfrak{a}_1'$ einer Klasse $\overline{A_1}$ an, und $\mathfrak{a}_2, \mathfrak{a}_2'$ einer Klasse
$\overline{A_2}$, so gehören die Produkte $\mathfrak{a}_1 \mathfrak{a}_2$ und $\mathfrak{a}_1' \mathfrak{a}_2'$ in \gesperrt{dieselbe Klasse}, die
wir die Klasse $\overline{A_1}\,\overline{A_2}$ nennen. Die Klassen lassen sich also komponieren und
bilden eine endliche \gesperrt{Abelsche Gruppe}, in der die
Hauptklasse $\overline{E}A$ die Einheit ist.

Eine Zahlgruppe $A$ soll \gesperrt{Kongruenzgruppe} heißen, wenn
sie nicht nur aus einzelnen Zahlen, sondern aus ganzen \gesperrt{Zahlklassen}
 nach einem Modul $M$ besteht, d.\,h. wenn $\alpha$ zu $A$ gehört,
so sollen auch alle mit $\alpha$ nach dem Modul $M$ kongruenten Zahlen
zu $A$ gehören; $M$ soll der \gesperrt{Modul der Kongruenzgruppe} heißen.
Eine Kongruenzgruppe hat immer einen endlichen Index, weil ja
schon die Anzahl der Zahlklassen nach dem Modul $M$ endlich ist.
Statt $M$ kann man auch jedes Vielfache von $M$ als Modul wählen.
Darum gewinnen wir auch keinen allgemeineren Begriff, wenn
wir für $M$ ein Ideal nehmen. Wir können sogar $M$ als natürliche
Zahl annehmen.

Wenn ein Ideal $\mathfrak{m}$ relativ prim zu $M$ ist, so kann man in
jeder durch $\alpha$ repräsentierten Zahlklasse eine Zahl $\alpha_0$ aus der
Kongruenz
%CONT>>
